\documentclass{article}
\usepackage{graphicx}
\usepackage{amsfonts}
\usepackage{amsmath}
\usepackage{amsthm}
\usepackage{tikz}

\theoremstyle{definition}
\newtheorem{theorem}{Theorem}[section]
\newtheorem{corollary}[theorem]{Corollary}
\newtheorem{proposition}[theorem]{Proposition}
\newtheorem{lemma}[theorem]{Lemma}
\newtheorem{definition}[theorem]{Definition}

\theoremstyle{remark}
\newtheorem{remark}[theorem]{Remark}
\newtheorem{example}[theorem]{Example}
\newtheorem{claim}[theorem]{Claim}

\title{Linear Transformation}
\author{Sarah Liu}
\date{January 2026}

\begin{document}

\maketitle


\section{Linear Transformation}
\begin{definition}
    A function $T:V\rightarrow W$ between vector space $V, W$ is a linear transformation if it satisfies conditions:
    \begin{enumerate}
        \item $T(\vec{x}+\vec{y})=T(\vec{x})+T(\vec{y})$
        \item $T(c\vec{x})=cT(\vec{x})$
    \end{enumerate}
\end{definition}

\begin{proposition}
    Let $T$ be a linear transformation, then $T(\vec{0_V})=\vec{0_W}$.
\end{proposition}

\begin{proof}
    For any $\vec{x} \in V$, $T(\vec{0_V})=T(0\cdot \vec{x})=0\cdot T(\vec{x})=\vec{0_W}$
\end{proof}

\begin{proposition}
    Let $T$ be a linear transformation and $S=\{\vec{v_1},\dots ,\vec{v_m}\}$ is a basis of $V$, then $T(\vec{v_1}), \dots ,T(\vec{v_m})$ determines $T$.
\end{proposition}

\begin{proof}
    For any $\vec{v}\in V$, we could express $\vec{v}$ uniquely as $\vec{v}=a_1\vec{v_1}+\dots+a_m\vec{v_m}$. Then 
    \[
    T(\vec{v})=a_1T(\vec{v_1})+ \dots +a_mT(\vec{v_m})
    \]
    so we conclude that $T(\vec{v_1}), \dots ,T(\vec{v_m})$ determines $T$.
\end{proof}

\section{Matrix Associated with Linear Transformation}

\begin{definition}[Coordinate Vector]
    Let $V$ be a vector space, and $\alpha=\{\vec{v_1},\dots, \vec{v_n}\}$ a basis of $V$. Then the coordinate for vector $\vec{v}\in V$ with respect to $\alpha$ is 
    \[
    [\vec{v}]_{\alpha}=
    \begin{pmatrix}
a_1 \\ \vdots \\ a_n
\end{pmatrix}
    \]
    where $\vec{v}=a_1\vec{v_1}+\dots+a_n\vec{v_n}$.
\end{definition}

\begin{definition}
    Let $T:V\rightarrow W$ be a linear transformation. $\alpha=\{\vec{v_1},\dots, \vec{v_n}\}$ is a basis of $V$, and $\beta = \{\vec{w_1},\dots, \vec{w_m}\}$ is a basis of $W$.
    We could write
    \[
    \begin{aligned}
        T(\vec{v_1})&= a_{11}\vec{w_1}+a_{21}\vec{w_2}+\dots+a_{m1}\vec{w_m}\\
        T(\vec{v_2})&= a_{12}\vec{w_1}+a_{22}\vec{w_2}+\dots+a_{m2}\vec{w_m}\\
        &\dots\\
        T(\vec{v_n})&= a_{1n}\vec{w_1}+a_{2n}\vec{w_2}+\dots+a_{mn}\vec{w_m}
    \end{aligned}
    \]
    Then the matrix of $T$ with respect to $\alpha, \beta$ is 
    \[
    [T]_{\alpha}^{\beta}=
    \begin{bmatrix}
    a_{11} & a_{12} & \dots & a_{1n} \\
    a_{21} & a_{22} & \dots & a_{2n} \\
    \dots & \dots & \dots & \dots\\
    a_{m1} & a_{m2} & \dots & a_{mn}
    \end{bmatrix}
    \]
\end{definition}

\begin{proposition}
    Let $T:V\rightarrow W$ be a linear transformation.
    $\alpha=\{\vec{v_1},\dots, \vec{v_n}\}$ is a basis of $V$, and $\beta = \{\vec{w_1},\dots, \vec{w_m}\}$ is a basis of $W$.
    Then we have
    \[
    [T]_{\alpha}^{\beta}[\vec{x}]_{\alpha}=[T(\vec{x})]_{\beta}
    \]
\end{proposition}

\begin{proof}
    Let $\vec{x}=a_1\vec{v_1}+\dots+a_n\vec{v_n}$, then $[\vec{x}]_{\alpha}=(a_1,\dots,a_n)$.
    Then,
    \[
    \begin{aligned}
        T(\vec{x})&=T(a_1\vec{v_1}+\dots+a_n\vec{v_n})\\
    &=a_1T(\vec{v_1})+ \dots +a_nT(\vec{v_n})\\
    &=a_1\sum_{j=1}^{m}a_{j1}\vec{w_j}+\dots+a_n\sum_{j=1}^{m}a_{jn}\vec{w_j}\\
    &=\vec{w_1}\sum_{i=1}^{n}a_ia_{1i}+\dots +\vec{w_m}\sum_{i=1}^{n}a_ia_{mi}
    \end{aligned}
    \]
    So, 
    \[
    [T(\vec{x})]_{\beta}=
    \begin{bmatrix}
        a_1a_{11}+\dots+a_na_{1n}\\
        \dots\\
        a_1a_{m1}+\dots+a_na_{mn}
    \end{bmatrix}
    \]
    Therefore,
    \[
    [T]_{\alpha}^{\beta}[\vec{x}]_{\alpha}=
    \begin{bmatrix}
    a_{11} & a_{12} & \dots & a_{1n} \\
    a_{21} & a_{22} & \dots & a_{2n} \\
    \dots & \dots & \dots & \dots\\
    a_{m1} & a_{m2} & \dots & a_{mn}
    \end{bmatrix}
    \begin{bmatrix}
        a_1\\ \vdots\\a_n
    \end{bmatrix}
    =
    \begin{bmatrix}
        a_1a_{11}+\dots+a_na_{1n}\\
        \dots\\
        a_1a_{m1}+\dots+a_na_{mn}
    \end{bmatrix}=[T(\vec{x})]_{\beta}
    \]
\end{proof}

\begin{example}
    The change of basis coordinate matrix is $[I]_{\alpha}^{\beta}$.
\end{example}

\section{Kernel, Image and Dimensional Theorem}

\begin{definition}[Kernel/Image]
    Let $T:V\rightarrow W$ be a linear transformation. 
    The \textbf{Kernel} of $T$ is the subset of $V$ that satisfies:
    \[
    \text{Ker}(T):=\{\vec{v}\in V:T(\vec{v})=\vec{0_W}\}
    \]
    The \textbf{Image} of $T$ is the subset of $W$ that satisfies:
    \[
    \text{Im}(T):=\{\vec{w}\in W:\exists \vec{v}\in V\;\text{s.t.}\;T(\vec{v})=\vec{w}\}
    \]
\end{definition}

\begin{theorem}
    Let $T:V\rightarrow W$ be a linear transformation. 
    \begin{enumerate}
        \item $\text{Ker}(T)$ is a subspace of $V$
        \item $\text{Im}(T)$ is a subspace of $W$
    \end{enumerate}
\end{theorem}

\begin{proof}
    We know from $T(\vec{0_V})=\vec{0_W}$ that $\vec{0_V}\in\text{Ker}(T)$.
    For $\vec{x},\vec{y}\in\text{Ker}(T)$, $T(\vec{x})=\vec{0_W}$, $T(\vec{y})=\vec{0_W}$.
    From the linearity of transformation $T$, we know that $T(\vec{x}+\vec{y})=T(\vec{x})+T(\vec{y})=\vec{0_W}+\vec{0_W}=\vec{0_W}$.
    Also, $T(c\vec{x})=cT(\vec{x})=c\vec{0_W}=\vec{0_W}$.
\end{proof}

\begin{theorem}
    Let $T:V\rightarrow W$ be a linear transformation between finite dimensional vector spaces.
    Let $\alpha$ be a basis for $V$, and $\beta$ be a basis for $W$. 
    Then we have
    \begin{enumerate}
        \item $\vec{x}\in\text{Ker}(T) \iff [\vec{x}]_{\alpha}\in\text{Nul}[T]_{\alpha}^{\beta}$
        \item $\vec{y}\in\text{Im}(T) \iff [\vec{y}]_{\beta}\in\text{Col}[T]_{\alpha}^{\beta}$
    \end{enumerate}
\end{theorem}

\begin{proof}
    Let $\vec{x}\in\text{Ker}(T)$, then $T(\vec{x})=\vec{0}$.
    Because $[T]_{\alpha}^{\beta}[\vec{x}]_{\alpha}=[T(\vec{x})]_{\beta}=\vec{0}$, so $[\vec{x}]_{\alpha}\in\text{Nul}[T]_{\alpha}^{\beta}$.

    Let $\vec{y}\in\text{Im}(T)$, then $\exists \vec{v}\in V \text{ s.t. }\vec{y}=T(\vec{v})$.
    Because $[\vec{y}]_{\beta}=[T(\vec{v})]_{\beta}=[T]_{\alpha}^{\beta}[\vec{v}]_{\alpha}$, so $[\vec{y}]_{\beta}\in\text{Col}[T]_{\alpha}^{\beta}$.
\end{proof}

\begin{theorem}
    Let $V$ be a vector space, and $\alpha=\{\vec{v_1},\dots, \vec{v_n}\}$ be a basis of $V$.
    Then $\{\vec{x_1},\dots, \vec{x_n}\}$ is linearly independent if and only if $\{[\vec{x_1}]_{\alpha},\dots, [\vec{x_n}]_{\alpha}\}$ is linearly independent.
\end{theorem}

\begin{proof}
$a_1\vec{x_1}+\dots+a_n\vec{x_n}=\vec{0}$ if and only if $a_1=\dots=a_n=0$.
    \[
    \begin{aligned}
        a_1\vec{x_1}+\dots+a_n\vec{x_n}\iff[a_1\vec{v_1}+\dots+a_n\vec{v_n}]_{\alpha}\iff a_1[\vec{v_1}]_{\alpha}+\dots+a_n[\vec{v_n}]_{\alpha}
    \end{aligned}
    \]
    so $a_1[\vec{v_1}]_{\alpha}+\dots+a_n[\vec{v_n}]_{\alpha}=[\vec{0}]_{\alpha}$ if and only if $a_1=\dots=a_n=0$.
    So $\{[\vec{x_1}]_{\alpha},\dots, [\vec{x_n}]_{\alpha}\}$ is linearly independent.
        
    The opposite could be reasoned using the same argument
\end{proof}

\begin{theorem}
    Let $T:V\rightarrow W$ be a linear transformation. 
    Let $V$ be a finite dimensional vector space. Then
    \[
    \dim (\text{Ker} (T))+\dim(\text{Im}(T))=\dim V
    \]
\end{theorem}

\begin{proof}
    Assume $W$ is finite dimensional.
    Let $\alpha$ be a basis of $V$, and $\beta$ be a basis of $W$. 
    Write $M=[T]_{\alpha}^{\beta}$. Then 
    \[
    \dim (\text{Ker} (T))+\dim(\text{Im}(T))=\dim (\text{Nul} )(M)+\dim(\text{Col}(M))= \text{cols}(M) =\dim V
    \]
\end{proof}

\begin{proof}[Alternative]
    Let $\{\vec{v_1},\vec{v_2},\dots,\vec{v_k}\}$ be a basis of $\ker(T)$.
    Extend the basis to $\{\vec{v_1},\vec{v_2},\dots,\vec{v_n}\}$ to be a basis of $V$.
    \begin{claim}
        $\{T(\vec{v_{k+1}}),\dots,T(\vec{v_n})\}$ is a basis of Im$(T)$
    \end{claim}
\begin{proof}
    Let $\vec{w}\in \text{Im}(T)$.
    Then there exists $\vec{v}\in V$, such that $T(\vec{v})=\vec{w}$.
    Write $\vec{v}$ as $\vec{v}=a_1\vec{v_1}+\dots+a_n\vec{v_n}$.
    From $\vec{v_1},\vec{v_2},\dots,\vec{v_k}\in\ker(T)$, we know that $T(\vec{v_1})=0,T(\vec{v_2})=0,\dots,T(\vec{v_k})=0$.
    Then
    \[
    \begin{aligned}
    \vec{w}=T(\vec{v})&=T(a_1\vec{v_1}+\dots+a_n\vec{v_n})\\
    &=a_1T(\vec{v_1})+\dots+a_nT(\vec{v_n})\\
    &=a_{k+1}T(\vec{v_{k+1}})+\dots+a_nT(\vec{v_n})\\
    &\in \text{Span}(T(\vec{v_{k+1}}),\dots,T(\vec{v_n}))
    \end{aligned}
    \]
    Next, we prove that $\{T(\vec{v_{k+1}}),\dots,T(\vec{v_n})\}$ is a linearly independent.
    \[
    \begin{aligned}
        &a_{k+1}T(\vec{v_{k+1}})+\dots+a_nT(\vec{v_n})=\vec{0}\\
        \iff &T(a_{k+1}\vec{v_{k+1}}+\dots+a_n\vec{v_n})=\vec{0}\\
        \iff & a_{k+1}\vec{v_{k+1}}+\dots+a_n\vec{v_n}=\vec{0}
    \end{aligned}
    \]
    Because $\{\vec{v_1},\dots,\vec{v_n}\}$ is linearly independent, so $a_{k+1}=\dots=a_n=0$
\end{proof}

\end{proof}

\section{Injectivity and Surjectivity}

\begin{definition}
    Let $T$ be a function
    \begin{enumerate}
        \item $T$ is injective (one-to-one) if $T(\vec{v_1})=T(\vec{v_2})\implies \vec{v_1}=\vec{v_2}$
        \item $T$ is surjective (onto) if for every $\vec{w}\in W$ there exists $\vec{v}\in V$ such that $T(\vec{v})=\vec{w}$.
        In other words, $W=\text{Im}(V)$.
    \end{enumerate}
\end{definition}

\begin{theorem}
    Let $T:V\rightarrow W$ be a linear transformation. 
    $T$ is injective if and only if $\text{Ker}(T)=\{\vec{0}\}$.
\end{theorem}

\begin{proof}
    Assume $T$ is injective. 
    For any $\vec{v}\in \text{Ker}(T)$, we have $T(\vec{v})=\vec{0_W}$.
    For every linear transformation, we have $T(\vec{0_V})=\vec{0_W}$.
    Since $T$ is injective, we have that $\vec{v}=\vec{0_V}$, so $\text{Ker}(T)=\{\vec{0}\}$.
    For the opposite direction, we assume $\text{Ker}(T)=\{\vec{0}\}$.
    Suppose $T(\vec{v_1})=T(\vec{v_2})$, then $T(\vec{v_1})-T(\vec{v_2})=0$, so $T(\vec{v_1}-\vec{v_2})=\vec{0_W}$.
    Because $T$ is injective, so $\vec{v_1}-\vec{v_2}=\vec{0_V}$, so $\vec{v_1}=\vec{v_2}$.
\end{proof}

\begin{theorem}
    Suppose $V$ and $W$ are finite dimensional vector spaces.
    \begin{enumerate}
        \item If $T$ is injective, then $\dim V\leq\dim W$.
        \item If $T$ is surjective, then $\dim W\leq\dim V$.
    \end{enumerate}
\end{theorem}

\begin{proof}
    Assume $T$ is injective. 
    Then $\dim(\text{Ker(T)})=0$. 
    So $\dim V=\dim(\text{Im}(T))\leq\dim W$.
    For the second statement, we assume that $T$ is surjective.
    Then $\dim(\text{Im}(T))=\dim W$, so $\dim V=\dim(\text{Ker(T)})+\dim(\text{Im}(T))\geq\dim(\text{Im}(T))=\dim W$
\end{proof}

\begin{corollary}
    Suppose $\dim V=\dim W$.
    $T$ is injective if and only if $T$ is surjective.
\end{corollary}

\begin{proof}
    Assume $T$ is injective, then $\dim W=\dim V=\dim(\text{Ker(T)})+\dim(\text{Im}(T))=\dim(\text{Im}(T))$.
    So $T$ is surjective.
    Assume $T$ is surjective, then $\dim(\text{Im}(T))=\dim W=\dim V=\dim(\text{Ker(T)})+\dim(\text{Im}$, so $\dim(\text{Ker(T)})=0$. 
    So $T$ is injective.
\end{proof}

\begin{definition}
    Let $T:V\rightarrow W$ be a linear transformation. 
    We say $T$ is an isomorphism if it is both injective and surjective.
    We say that $V$ is isomorphic to $W$.
\end{definition}

\section{Composition}

\begin{definition}
    Let  $U,V,W$ be vector spaces. And $S:U\rightarrow V$, $T:V\rightarrow W$ the linear transformation. The composition of $S$ and $T$ is the function $TS:U\rightarrow W$ defined by $TS(\vec{u})=T(S(\vec{u}))$.
\end{definition}

\begin{proposition}
    Let $S:U\rightarrow V$, $T:V\rightarrow W$ be linear transformations. Then $TS:U\rightarrow W$ is also a linear transformation.
\end{proposition}

\begin{proof}
    Let $\vec{x},\vec{y}\in U,c\in\mathbb{R}$.
    Then $TS(\vec{x}+\vec{y})=T(S(\vec{x}+\vec{y}))=T(S(\vec{x})+S(\vec{y}))=T(S(\vec{x}))+T(S(\vec{y}))=TS(\vec{x})+TS(\vec{y})$.
    $TS(c\vec{x})=T(S(c\vec{x}))=T(cS(\vec{x}))=cT(S(\vec{x}))=cTS(\vec{x})$.
\end{proof}

\begin{theorem}
    Let $S:U\rightarrow V$, $T:V\rightarrow W$ be linear transformations on finite dimensional vector spaces.
    $\alpha,\beta,\gamma$ are basis of $U,V,W$ respectively. 
    Then
    \[
    [TS]_{\alpha}^{\gamma}=[T]_{\beta}^{\gamma}[S]_{\alpha}^{\beta}
    \]
\end{theorem}

\begin{proof}
    $[TS]_{\alpha}^{\gamma}[\vec{x}]_{\alpha}=[TS(\vec{x})]_{\gamma}=[T(S(\vec{x}))]_{\gamma}=[T]_{\beta}^{\gamma}[S(\vec{x})]_{\beta}=[T]_{\beta}^{\gamma}[S]_{\alpha}^{\beta}[\vec{x}]_{\alpha}$
\end{proof}

\section{Inverse}

\begin{definition}
    Let $T:V\rightarrow W$ a linear transformation. 
    An inverse $S:W\rightarrow V$ is a function, satisfies 
    \[
    \begin{aligned}
        ST(\vec{v})&=\vec{v}\quad \text{for all}\;\vec{v}\in V\\
        TS(\vec{w})&=\vec{w}\quad \text{for all}\;\vec{w}\in W
    \end{aligned}
    \]
\end{definition}

\begin{proposition}
    Let $T:V\rightarrow W$ a linear transformation. 
    Let $S:W\rightarrow V$ be the inverse of $T$.
    Then $S$ is linear.
\end{proposition}

\begin{proof}
    For $\vec{z},\vec{w}\in W$.
    Let $S(\vec{z})=\vec{x}$, then $T(\vec{x})=\vec{z}$.
    Let $S(\vec{w})=\vec{y}$, then $T(\vec{x})=\vec{w}$.
    Then
    \[
    \begin{aligned}
        &S(\vec{z}+\vec{w})=S(T(\vec{x})+T(\vec{y}))=S(T(\vec{x}+\vec{y}))=\vec{x}+\vec{y}=S(\vec{z})+S(\vec{w})\\
        &S(c\vec{z})=S(cT(\vec{x}))=S(T(c\vec{x}))=c\vec{x}=cS(\vec{z})
    \end{aligned}
    \]
\end{proof}

\begin{proposition}
    $T:V\rightarrow W$ has an inverse if and only if $T$ is isomorphism.
\end{proposition}

\begin{proposition}
    Let $T:V\rightarrow W$ be a isomorphism between finite dimensional vector spaces. Then $\dim (V)=\dim (W)$.
\end{proposition}

\begin{proof}
$\dim (V)=\dim (\text{Ker}(T))+\dim (\text{Im}(T))=\dim (W)$
\end{proof}

\begin{theorem}
    Let $V,W$ be finite dimensional vector spaces.
    If $\dim (V)=\dim(W)$, then there exists an isomorphism $T:V\rightarrow W$.
\end{theorem}

\begin{proof}
    Let $\alpha=\{\vec{v_1},\dots, \vec{v_n}\}$ be a basis of $V$, and $\beta=\{\vec{w_1},\dots, \vec{w_n}\}$ be a basis of $W$.
    Define $T(\vec{v_j})=\vec{w_j}$ for all $j=1,\dots,n$.

    First, prove that $T$ is injective.
    Write $\vec{v}=a_1\vec{v_1}+\dots+a_n\vec{v_n}$.
    Then $\vec{0_W}=T(\vec{v})=T(a_1\vec{v_1}+\dots+a_n\vec{v_n})=a_1T(\vec{v_1})+\dots+a_nT(\vec{v_n})=a_1\vec{w_1}+\dots+a_n\vec{w_n}$.
    Because $\{\vec{w_1},\dots, \vec{w_n}\}$ is a basis, so $a_1=\dots =a_n=0$, so $\vec{v}=\vec{0}$.

    Next, prove that $T$ is surjective.
    For $\vec{w}\in W$, write $\vec{w}=b_1\vec{w_1}+\dots+b_n\vec{w_n}$.
    Then $T(b_1\vec{v_1}+\dots+b_n\vec{v_n})=b_1T(\vec{v_1})+\dots+b_nT(\vec{v_n})=b_1\vec{w_1}+\dots+b_n\vec{w_n}=\vec{w}$, so the image of $T$ is $W$.

\end{proof}

\begin{proposition}
    Let $T:V\rightarrow W$ be a isomorphism between finite dimensional vector spaces.
    $T^{-1}$ is the inverse of $T$.
    Let $\alpha,\beta$ be basis of $V,W$.
    Then 
    \[
    [T^{-1}]_{\beta}^{\alpha}=([T]_{\alpha}^{\beta})^{-1}
    \]
\end{proposition}

\begin{proof}
    $[T^{-1}]_{\beta}^{\alpha}[T]_{\alpha}^{\beta}=[T^{-1}T]_{\alpha}^{\alpha}=[I]_{\alpha}^{\alpha}=I$
\end{proof}

\end{document}